\subsection{Feasibility Indicators}
This group reflects the concerns of any skeptical party that might argue about the feasibility of the commissioning process itself. These indicators capture all practical factors that may influence implementation ease, with evaluations conducted using qualitative scales based on current conditions (status quo) as detailed in the following analysis. All of the indicators in this group are qualitative in nature, the detailed methodology to charaterize them is presented in section \ref{sec:qi}. Hereafter, each indicator is described as well as the reason behind their starting point evaluation.

\paragraph{Design Maturity}
Design maturity was evaluated on an ordinal scale reflecting current deployment status. Among large reactors, the \gls{epr} received a 4/5 score due to its limited construction experience, while both \gls{apr} and \gls{ap} achieved 5/5 ratings. For \gls{smr} designs, only the Rolls-Royce reactor scored 2/5 rather than 1/5, as its design closely resembles existing \glspl{pwr} featuring forced circulation and external steam generators. Table \ref{tab:maturity-comparison} presents the supporting data for this preliminary maturity assessment.

\begin{table}[htbp]
    \centering
    \renewcommand{\tabularxcolumn}[1]{m{#1}}
    \begin{tabularx}{\textwidth}{Xcccccc}
        \toprule
        \textbf{Parameter} & \textbf{APR} & \textbf{EPR} & \textbf{AP} & \textbf{BWRX} & \textbf{Rolls Royce} & \textbf{NuScale} \\
        \midrule
        Operational units & 9 & 2 (+2) & 7 & 0 & 0 & 0 \\
        \cmidrule(lr){1-7}
        Under construction & 3 & 2 & 2 & 0 & 0 & 0 \\
        \cmidrule(lr){1-7}
        Technology & PWR & PWR & PWR & i-BWR & PWR & i-PWR \\
        \cmidrule(lr){1-7}
        Construction strategy (Units) & 1 & 1 & 1 & 1 & 1 & 4–12  \\
        \cmidrule(lr){1-7}
        Fuel & Existing & Existing & Existing & Existing & Variation & Variation \\
        \cmidrule(lr){1-7}
        Design Maturity Score & 5/5 & 4/5 & 5/5 & 1/5 & 2/5 & 1/5 \\
        \bottomrule
    \end{tabularx}
    \caption{Data to support design maturity.}
    \label{tab:maturity-comparison}
\end{table}


\paragraph{Licensing Status}
Licensing progress was evaluated based on each reactor current regulatory status within the European context. The three conventional reactors received baseline scores of 4/5 due to their established international licensing history, with the \gls{epr} achieving 5/5 given its existing European deployments, with the previous iteration, in France and Finland \footnote{For clarity regarding nuclear reactor nomenclature: Internet searches for "EPR" will present results for Flamanville and Olkiluoto plants as well as UK units (Hinkley Point and Sizewell) and China's Taishan plant. This is actually an aggregation of two distinct designs: the original 1600 MWe "European Pressurized Reactor" (Flamanville/Olkiluoto) and the evolved 1750 MWe "Evolutionary Pressurized Reactor" (Hinkley Point/Sizewell/Taishan) analyzed in this study.}.


For \glspl{smr} designs, NuScale and Rolls-Royce received 1/5 scores as their licensing processes remain confined to overseas markets (United States for NuScale; United States and United Kingdom for Rolls-Royce). The \gls{bwrx} was assigned 2/5 despite similar overseas licensing status (United States, United Kingdom, and Canada) to acknowledge three distinguishing factors: exceptional financing support (scoring 6/6 in the \gls{nea} \gls{smr} Dashboard \cite{NEA_SMR_Dashboard_2025}), the issue of a Canadian building permit, and utilization of existing fuel designs that eliminate requirements for new fuel licensing procedures. The 3/5 score was intentionally omitted to clearly differentiate between established and emerging reactor designs. Table \ref{tab:licensing_info} summarizes the licensing status evaluation.
\begin{table}[!hb]
    \begin{tabularx}{\textwidth}{l>{\centering\arraybackslash}X}
        \toprule
        \textbf{Reactor Model} & \textbf{Information} \\
        \midrule
        APR-1400 &
          \begin{itemize}
            \item Built in Korea, UAE
            \item Licensed in US
            \item Approved by EUR
          \end{itemize} \\
        \midrule
        EPR-1750 &
          \begin{itemize}
            \item Built in UK, China (+ FR, FIN)
            \item Approved by EUR
          \end{itemize} \\
        \midrule
        AP-1000 &
          \begin{itemize}
            \item Built in China, US
            \item Approved by EUR
          \end{itemize} \\
        \midrule
        BWRX-300 &
          \begin{itemize}
            \item Build license in CAN 2025 (subject to VDR 2023)
            \item Pre-application in US in 2019
            \item UK GDA started 2024
          \end{itemize} \\
        \midrule
        Rolls Royce SMR &
          \begin{itemize}
            \item Pre-application in US in 2025
            \item UK GDA started in 2022
          \end{itemize} \\
        \midrule
        NuScale &
          \begin{itemize}
            \item SDA in US (2025)
          \end{itemize} \\
        \bottomrule
    \end{tabularx}
    \caption{Licensing and construction information for selected reactor models.}
    \label{tab:licensing_info}
\end{table}

\paragraph{Supplier Availability}
Market entry ease and local industry involvement significantly affect commissioning timelines and national economic benefits, both of which are key considerations for governmental agencies. While a comprehensive evaluation would require detailed analysis of each country's nuclear industry status and expansion potential, this preliminary study implements a simplified assessment approach.

The methodology employs a country-reactor matrix (with each reactor representing a different vendor) to evaluate two primary factors: supplier presence and potential industry capacity. The initial scoring assigns +0.25 points for supplier legal presence through registered offices and an additional +0.25 points for existing fabrication capabilities. These allocations produce the baseline scores presented in Table \ref{tab:supplier-availability-step1}.

\begin{table}[!ht]
    \centering
    \begin{tabularx}{\textwidth}{X c c c c c c}
        \toprule
        \textbf{Country} & \textbf{APR} & \textbf{EPR} & \textbf{AP} & \textbf{BWRX} & \textbf{RR SMR} & \textbf{NuScale} \\
        \midrule
        France (FR)       & 0.00 & 0.50 & 0.50 & 0.25 & 0.25 & 0.00 \\
        Italy (IT)       & 0.00 & 0.25 & 0.25 & 0.25 & 0.00 & 0.00 \\
        Switzerland (CH)  & 0.00 & 0.25 & 0.25 & 0.25 & 0.00 & 0.00 \\
        Poland (PL)       & 0.00 & 0.00 & 0.25 & 0.25 & 0.00 & 0.00 \\
        \bottomrule
    \end{tabularx}
    \caption{First step of supplier availability scores for each reactor model in the selected countries.}
    \label{tab:supplier-availability-step1}
\end{table}

The broader nuclear industry is also evaluated to consider domestic firms qualified to manufacture nuclear-certified materials and components, essential for the "nuclear island". Nuclear certification requirements extend beyond basic materials to include specialized components: structural steel for pressure vessels must conform to stringent nuclear standards, primary coolant pumps require dedicated certification, and certain elements demand highly specialized fabrication techniques\footnote{For instance, \gls{epr} reactor pressure vessels require precision manufacturing of a 4.8m diameter, 13m height steel cylinder with walls 25cm thick.}. If the country presents any level of industry to support nuclear activites gets a +0.25, if it also has capability of manufacturing large \gls{nsss} components it get an additional +0.25. For this reason Switerland get a +0.25 while both France and Italy a +0.5. The combined scores from both evaluation phases are presented in Table \ref{tab:supplier-availability-step2}.

\begin{table}[!ht]
    \centering
    \begin{tabularx}{\textwidth}{X c c c c c c}
        \toprule
        \textbf{Country} & \textbf{APR} & \textbf{EPR} & \textbf{AP} & \textbf{BWRX} & \textbf{RR SMR} & \textbf{NuScale} \\
        \midrule
        France (FR)       & 0.50 & 1.00 & 1.00 & 0.75 & 0.75 & 0.50 \\
        Italy (IT)       & 0.50 & 0.75 & 0.75 & 0.75 & 0.50 & 0.50 \\
        Switzerland (CH)  & 0.25 & 0.50 & 0.50 & 0.50 & 0.25 & 0.25 \\
        Poland (PL)       & 0.00 & 0.00 & 0.25 & 0.25 & 0.00 & 0.00 \\
        \bottomrule
    \end{tabularx}
    \caption{Final supplier availability scores for each reactor model in the selected countries.}
    \label{tab:supplier-availability-step2}
\end{table}

Two significant considerations emerge from this analysis. First, while Poland is currently experiencing a period of strong economic growth, the potential acceleration of its nuclear sector remains uncertain and challenging to quantify. Second, established supply chains were excluded from France's evaluation due to the country having completed only one nuclear power plant (Flamanville-3) in the past 25 years, likely indicating supply chain degradation despite its historical capabilities in this sector.

